% Однофазный мост VS1-VS4 с индуктивностью рассеяния Ls (эквивалент
% по Коновалову рис. 5.4: одна обмотка, одна Ls — рассеяние всей
% вторичной обмотки). Тиристоры — прямоугольник; интервал коммутации
% gamma — открыты все четыре
\documentclass[border=4pt]{standalone}
\usepackage{amsmath}
\usepackage{fontspec}
\setmainfont{Calibri}
\usepackage[european, RPvoltages]{circuitikz}
\definecolor{ink}{RGB}{31,42,51}
\definecolor{grp}{RGB}{46,110,140}
\ctikzset{bipoles/length=1.0cm}
\ctikzset{bipoles/thickness=0.9}
\ctikzset{diodes/fill=ink}
\ctikzset{inductor=american}

\begin{document}
\begin{tikzpicture}[line width=1.0pt, color=ink, font=\small]
  % --- вторичная обмотка (дуги полуколец — влево)
  \draw (0,3.0) to[L] (0,4.2);
  % --- верхний вывод: единственная Ls (рассеяние всей обмотки)
  %     до средней точки левой ноги
  \draw (0,4.2) -- (0.6,4.2) to[L, l=$L_{\text{s}}$] (2.0,4.2) -- (2.6,4.2) node[circ]{};
  % --- нижний вывод: провод к средней точке правой ноги
  %     (перекрёсток с левой ногой в (2.6,3.0) — без точки)
  \draw (0,3.0) -- (0.6,3.0) -- (4.6,3.0) node[circ]{};
  % --- левая нога: VS3 (внизу) и VS1 (вверху) на одной вертикали
  \draw (2.6,0.6) -- (2.6,1.4)
        to[Ty, name=VS3] (2.6,2.2) -- (2.6,4.9)
        to[Ty, name=VS1] (2.6,5.7) -- (2.6,6.0);
  % --- правая нога: VS4 (внизу) и VS2 (вверху) на одной вертикали
  \draw (4.6,0.6) node[circ]{} -- (4.6,1.4)
        to[Ty, name=VS4] (4.6,2.2) -- (4.6,3.0) node[circ]{} -- (4.6,4.9)
        to[Ty, name=VS2] (4.6,5.7) -- (4.6,6.0) node[circ]{};
  % --- шины + и - ; нагрузка: дроссель Ld + Rн
  \draw (2.6,6.0) -- (7.0,6.0);
  \draw (2.6,0.6) -- (7.0,0.6);
  \draw (7.0,6.0) to[L, l=$L_{\text{d}}$] (7.0,3.9)
        -- (7.0,3.1) to[R, l=$R_{\text{н}}$] (7.0,1.3) -- (7.0,0.6);
  % --- подписи
  \node[anchor=west] at (2.95,5.3) {VS1};
  \node[anchor=west] at (4.95,5.3) {VS2};
  \node[anchor=west] at (2.95,1.8) {VS3};
  \node[anchor=west] at (4.95,1.8) {VS4};
  \node[anchor=east] at (6.6,4.6) {$i_{\text{d}}$};
  \node[anchor=south west] at (6.95,6.02) {$+$};
  \node[anchor=north west] at (6.95,0.58) {$-$};
  % --- подпись группы
  \node[anchor=south, text=grp, font=\footnotesize] at (3.6,6.4)
       {интервал коммутации $\gamma$: открыты ВСЕ четыре тиристора, $u_{\text{d}} = 0$};
\end{tikzpicture}
\end{document}
